% Nouveau EGA V (cote 25): archive cote 25, site page 134 (PDF page 135).
% typescript. Draft compared: none.
% AI-transcribed by Claude Opus 5.5 (claude-opus-5-5, adaptive thinking, effort high), Message Batches API, batch ega25b (msgbatch_01A7n5YuXubUPNK1zsepi8RG), 2026-11-13; not yet checked by eye.
\scanpage{25/134}
\mspage{\hand{37.}}

\margin{29.9.1965}

Cher Dieudonné,

Merci de ta lettre du 24  et pour la table des matières des par.
16 à 19.  Je serais content de recevoir à l'occasion la table des
matières provisoire des par.20 et 21 ;\xout{\unsure{d}}d'accord pour les joindre au
fascicule 4 du Chap IV. Mais comment vaux-tu \xout{les} subdiviser\inserted{mon ancien par.20}, et
quels seront les titres des deux morceaux \unsure{?|2} Comme je commence à me
perdre dans le plan, et qu'il est \unsure{paffois} commode de \unsure{pauvoir|pouvoir} référer
sans trop déconner à un nº de paragraphe, je te donne ici ce qui me
semble être le plan actuel, dis-moi si tu es d'accord:

20.  ???

21.  ???

22.  Systêmes linéaires, compléments sur le groupe de Picard.

23.  Grassmaniennes.

24.  Formes lisses, singularités quadratiques ordinaires.

25.  Sections hyperplanes et bordel.

26.  Résultant et discriminant.

27.  Extensions infinitésimales.

Le 25.  risque d'\unsure{a|q}illeurs d'être fort long, et je te vois déjà
vouloir le subdiviser en d\unsure{eux} !  Pourtant, $27=3^3$ est un bien joli
nombre !

Il n'est pas question que je publie l'ex-Appendice au par.18
sous mon nom; ta rédaction n'a à peu près plus rien de commun avec
les vagues notes manuscriptes que je t'avais passées, si même je t'en
ai jamais passé, et ne me suis borné à te dire: il n'y a qu'à faire

\mspage{\hand{134}}
